\scanpage{60}
% archive foliation: 59

\textcircled{f}
\[
\begin{aligned}
X,Y\in\mathcal{E}_{\chi!}&\Longrightarrow X\times Y\in\mathcal{E}_{\chi!}\quad\text{et on a}\quad\chi_!(X\times Y)=\chi_!(X)\chi_!(Y)\\
X,Y\in\mathcal{E}_{\chi}&\Longrightarrow X\times Y\in\mathcal{E}_{\chi}\quad\text{et on a}\quad\chi(X\times Y)=\chi(X)\chi(Y).
\end{aligned}
\]
(La première assertion résulte de Künneth pour $H^*_!(X\times Y,k)$, la deuxième de Künneth pour $H_*(X\times Y,k)$, dont le dual est $H^*(X\times Y,k)$, donc qui permet de tester l'appartenance à $\mathcal{E}_{\chi}$ et la valeur de $\chi$ ...)

\textcircled{g} Soit $f:X\unsure{-|\longrightarrow}Y$ fibration loc. triviale de fibre $F$.

Si $Y,F\in\mathcal{E}_{\chi!}$, alors $X\in\mathcal{E}_{\chi!}$, et on a
\[
\chi_!(X)=\chi_!(Y)\chi_!(F).
\]
Énoncé analogue pour \inserted{$Y,F\in$}$\mathcal{E}_{\chi}$ et \deleted{$\unsure{\chi(X)=}$}
\[
\chi(X)=\chi(Y)\chi(F)
\]
\margin{\unsure{cum|avec} \unsure{léger} grain de sel}

\uline{Dém}. On a
\[
H^*_!(X)\Longleftarrow E_2^{pq}=H^p_!\left(Y,\underbrace{R^q f_!(\unsure{k_X})}_{\substack{\text{syst. local sur }Y\\\text{à fibres }\sim H^q_!(F)}}\right)
\]
donc \uline{sous réserve} qu'on sache que $E_2^{\unsure{**|pq}}$ de dim finie, on trouve $X\in\mathcal{E}_{\chi}$ et
\[
\begin{aligned}
\chi(X)&=\sum_{pq}(-1)^{p+q}\operatorname{rg}_k H^p_!(Y,R^q f_!(k_X))\\
&=\sum_q(-1)^q\underbrace{\sum_p(-1)^p\operatorname{rg}_k H^p_!(Y,R^q f_!(k_X))}_{b^q_!(F)\chi_!(Y)}\\
&=\chi_!(Y)\deleted{\unsure{\chi}}\sum_q(+1)^q b^q_!(F)=\chi_!(Y)\chi_!(F)
\end{aligned}
\]

Il faut prouver

\uline{Lemme} Soit $\mathcal{R}$ syst. local de $k$-vectoriels sur \deleted{\unsure{$X$ loc cpt,}} $Y\in\mathcal{E}_{\chi!}$, à fibres de rang \inserted{fini} $b$. Alors
\[
\chi_!(Y,\mathcal{R})\overset{\mathrm{déf}}=\sum_p(-1)^p H^p(Y,\mathcal{R})
\]
existe (i.e. $\unsure{\sum|\bigoplus}_p H^p(Y,\mathcal{R})$ de dim finie sur $k$) et on a
\[
\chi_!(Y,\mathcal{R})=b\chi_!(Y).
\]
\margin{\textcircled{\unsure{gn}}}

Marche tout au moins si $Y$ \deleted{triangulable} \inserted{admet filtration finie} avec $Y_i-Y_{i-1}$ somme finie de cellules \inserted{$Z_{ij}$}$\simeq R^{\delta_{ij}}$, car

On a :
\[
\begin{aligned}
\chi_!(Y,\mathcal{R})&=\sum_{ij}\chi_!(Z_{ij},\mathcal{R})=\sum_{ij}b\chi_!(Z_{ij})\\
&=b\sum\chi_!(Z_{ij})=b\chi_!(Y).
\end{aligned}
\]
\margin{(\unsure{ct.} sur $Z_{ij}$\unsure{)|]}}

